\documentclass[12pt,a4paper]{exam}
\usepackage[utf8]{inputenc}
\usepackage{amsmath,amssymb,amsthm}
\usepackage[margin=1in]{geometry}
\usepackage{graphicx}
\usepackage[colorlinks=true]{hyperref}
\pagestyle{headandfoot}
\firstpageheader{MIT OpenCourseWare}{}{\textbf{EXTREMELY HARD -- MIT LEVEL}}
\runningheader{MIT OpenCourseWare}{}{\textbf{EXTREMELY HARD -- MIT LEVEL}}
\firstpagefooter{}{Page \thepage}{}
\runningfooter{}{Page \thepage}{}

\begin{document}

\begin{center}
{\LARGE \textbf{MIT OpenCourseWare}}\\7.28\\Molecular Biology
\vspace{0.5cm}
{7.28} --- {Molecular Biology}
\vspace{0.3cm}
May 2026
\vspace{0.3cm}
\textbf{EXTREMELY HARD -- MIT LEVEL}
\end{center}

\begin{center}
\textbf{Time Limit: 3 hours}
\end{center}

\begin{center}
\textbf{Total Marks: 100}
\end{center}

\vspace{0.5cm}

\begin{questions}

\question[2]
Prove the fundamental theorem concerning DNA Replication in the context of 7.28 (Molecular Biology). State the theorem precisely, provide a complete proof, and discuss three important corollaries.

\vspace{0.3cm}

\question[4]
Prove the fundamental theorem concerning Transcription in the context of 7.28 (Molecular Biology). State the theorem precisely, provide a complete proof, and discuss three important corollaries.

\vspace{0.3cm}

\question[2]
Analyze the limitations of standard approaches to RNA Processing when applied to edge cases in 7.28. Construct explicit counterexamples showing where intuitive reasoning fails.

\vspace{0.3cm}

\question[8]
Prove the fundamental theorem concerning Translation in the context of 7.28 (Molecular Biology). State the theorem precisely, provide a complete proof, and discuss three important corollaries.

\vspace{0.3cm}

\question[3]
Construct an original proof-level problem involving Gene Regulation for 7.28 (Molecular Biology) and provide a complete solution. Your problem should be at the level of a graduate qualifying exam.

\vspace{0.3cm}

\question[3]
Construct an original proof-level problem involving Chromatin for 7.28 (Molecular Biology) and provide a complete solution. Your problem should be at the level of a graduate qualifying exam.

\vspace{0.3cm}

\question[3]
Construct an original proof-level problem involving DNA Replication for 7.28 (Molecular Biology) and provide a complete solution. Your problem should be at the level of a graduate qualifying exam.

\vspace{0.3cm}

\question[6]
Prove the fundamental theorem concerning Transcription in the context of 7.28 (Molecular Biology). State the theorem precisely, provide a complete proof, and discuss three important corollaries.

\vspace{0.3cm}

\question[5]
Analyze the limitations of standard approaches to RNA Processing when applied to edge cases in 7.28. Construct explicit counterexamples showing where intuitive reasoning fails.

\vspace{0.3cm}

\question[5]
Construct an original proof-level problem involving Translation for 7.28 (Molecular Biology) and provide a complete solution. Your problem should be at the level of a graduate qualifying exam.

\vspace{0.3cm}

\question[4]
Analyze the limitations of standard approaches to Gene Regulation when applied to edge cases in 7.28. Construct explicit counterexamples showing where intuitive reasoning fails.

\vspace{0.3cm}

\question[3]
Analyze the limitations of standard approaches to Chromatin when applied to edge cases in 7.28. Construct explicit counterexamples showing where intuitive reasoning fails.

\vspace{0.3cm}

\question[4]
Construct an original proof-level problem involving DNA Replication for 7.28 (Molecular Biology) and provide a complete solution. Your problem should be at the level of a graduate qualifying exam.

\vspace{0.3cm}

\question[7]
Analyze the limitations of standard approaches to Transcription when applied to edge cases in 7.28. Construct explicit counterexamples showing where intuitive reasoning fails.

\vspace{0.3cm}

\question[3]
Prove the fundamental theorem concerning RNA Processing in the context of 7.28 (Molecular Biology). State the theorem precisely, provide a complete proof, and discuss three important corollaries.

\vspace{0.3cm}

\question[5]
Analyze the limitations of standard approaches to Translation when applied to edge cases in 7.28. Construct explicit counterexamples showing where intuitive reasoning fails.

\vspace{0.3cm}

\question[3]
Construct an original proof-level problem involving Gene Regulation for 7.28 (Molecular Biology) and provide a complete solution. Your problem should be at the level of a graduate qualifying exam.

\vspace{0.3cm}

\question[3]
Analyze the limitations of standard approaches to Chromatin when applied to edge cases in 7.28. Construct explicit counterexamples showing where intuitive reasoning fails.

\vspace{0.3cm}

\question[3]
Construct an original proof-level problem involving DNA Replication for 7.28 (Molecular Biology) and provide a complete solution. Your problem should be at the level of a graduate qualifying exam.

\vspace{0.3cm}

\question[5]
Analyze the limitations of standard approaches to Transcription when applied to edge cases in 7.28. Construct explicit counterexamples showing where intuitive reasoning fails.

\vspace{0.3cm}

\question[4]
Analyze the limitations of standard approaches to RNA Processing when applied to edge cases in 7.28. Construct explicit counterexamples showing where intuitive reasoning fails.

\vspace{0.3cm}

\question[5]
Construct an original proof-level problem involving Translation for 7.28 (Molecular Biology) and provide a complete solution. Your problem should be at the level of a graduate qualifying exam.

\vspace{0.3cm}

\question[3]
Prove the fundamental theorem concerning Gene Regulation in the context of 7.28 (Molecular Biology). State the theorem precisely, provide a complete proof, and discuss three important corollaries.

\vspace{0.3cm}

\question[5]
Prove the fundamental theorem concerning Chromatin in the context of 7.28 (Molecular Biology). State the theorem precisely, provide a complete proof, and discuss three important corollaries.

\vspace{0.3cm}

\question[2]
Prove the fundamental theorem concerning DNA Replication in the context of 7.28 (Molecular Biology). State the theorem precisely, provide a complete proof, and discuss three important corollaries.

\vspace{0.3cm}

\end{questions}

\newpage
\section*{Solutions}

\textbf{Solution 1:}

The proof follows from the definitions and properties established in 7.28. The key insight is to apply the relevant theorem and verify the hypotheses hold. Detailed solution depends on the specific formulation of the theorem.

\vspace{0.5cm}

\textbf{Solution 2:}

The proof follows from the definitions and properties established in 7.28. The key insight is to apply the relevant theorem and verify the hypotheses hold. Detailed solution depends on the specific formulation of the theorem.

\vspace{0.5cm}

\textbf{Solution 3:}

The edge case analysis reveals the subtle assumptions underlying standard treatments of RNA Processing. By constructing explicit counterexamples, we see that the usual hypotheses cannot be weakened without loss of the theorem's conclusion.

\vspace{0.5cm}

\textbf{Solution 4:}

The proof follows from the definitions and properties established in 7.28. The key insight is to apply the relevant theorem and verify the hypotheses hold. Detailed solution depends on the specific formulation of the theorem.

\vspace{0.5cm}

\textbf{Solution 5:}

Solution approach: First recall the definitions and key theorems related to Gene Regulation from 7.28. Then apply the appropriate techniques to construct a rigorous argument. The solution must be self-contained and reference only the material from the course.

\vspace{0.5cm}

\textbf{Solution 6:}

Solution approach: First recall the definitions and key theorems related to Chromatin from 7.28. Then apply the appropriate techniques to construct a rigorous argument. The solution must be self-contained and reference only the material from the course.

\vspace{0.5cm}

\textbf{Solution 7:}

Solution approach: First recall the definitions and key theorems related to DNA Replication from 7.28. Then apply the appropriate techniques to construct a rigorous argument. The solution must be self-contained and reference only the material from the course.

\vspace{0.5cm}

\textbf{Solution 8:}

The proof follows from the definitions and properties established in 7.28. The key insight is to apply the relevant theorem and verify the hypotheses hold. Detailed solution depends on the specific formulation of the theorem.

\vspace{0.5cm}

\textbf{Solution 9:}

The edge case analysis reveals the subtle assumptions underlying standard treatments of RNA Processing. By constructing explicit counterexamples, we see that the usual hypotheses cannot be weakened without loss of the theorem's conclusion.

\vspace{0.5cm}

\textbf{Solution 10:}

Solution approach: First recall the definitions and key theorems related to Translation from 7.28. Then apply the appropriate techniques to construct a rigorous argument. The solution must be self-contained and reference only the material from the course.

\vspace{0.5cm}

\textbf{Solution 11:}

The edge case analysis reveals the subtle assumptions underlying standard treatments of Gene Regulation. By constructing explicit counterexamples, we see that the usual hypotheses cannot be weakened without loss of the theorem's conclusion.

\vspace{0.5cm}

\textbf{Solution 12:}

The edge case analysis reveals the subtle assumptions underlying standard treatments of Chromatin. By constructing explicit counterexamples, we see that the usual hypotheses cannot be weakened without loss of the theorem's conclusion.

\vspace{0.5cm}

\textbf{Solution 13:}

Solution approach: First recall the definitions and key theorems related to DNA Replication from 7.28. Then apply the appropriate techniques to construct a rigorous argument. The solution must be self-contained and reference only the material from the course.

\vspace{0.5cm}

\textbf{Solution 14:}

The edge case analysis reveals the subtle assumptions underlying standard treatments of Transcription. By constructing explicit counterexamples, we see that the usual hypotheses cannot be weakened without loss of the theorem's conclusion.

\vspace{0.5cm}

\textbf{Solution 15:}

The proof follows from the definitions and properties established in 7.28. The key insight is to apply the relevant theorem and verify the hypotheses hold. Detailed solution depends on the specific formulation of the theorem.

\vspace{0.5cm}

\textbf{Solution 16:}

The edge case analysis reveals the subtle assumptions underlying standard treatments of Translation. By constructing explicit counterexamples, we see that the usual hypotheses cannot be weakened without loss of the theorem's conclusion.

\vspace{0.5cm}

\textbf{Solution 17:}

Solution approach: First recall the definitions and key theorems related to Gene Regulation from 7.28. Then apply the appropriate techniques to construct a rigorous argument. The solution must be self-contained and reference only the material from the course.

\vspace{0.5cm}

\textbf{Solution 18:}

The edge case analysis reveals the subtle assumptions underlying standard treatments of Chromatin. By constructing explicit counterexamples, we see that the usual hypotheses cannot be weakened without loss of the theorem's conclusion.

\vspace{0.5cm}

\textbf{Solution 19:}

Solution approach: First recall the definitions and key theorems related to DNA Replication from 7.28. Then apply the appropriate techniques to construct a rigorous argument. The solution must be self-contained and reference only the material from the course.

\vspace{0.5cm}

\textbf{Solution 20:}

The edge case analysis reveals the subtle assumptions underlying standard treatments of Transcription. By constructing explicit counterexamples, we see that the usual hypotheses cannot be weakened without loss of the theorem's conclusion.

\vspace{0.5cm}

\textbf{Solution 21:}

The edge case analysis reveals the subtle assumptions underlying standard treatments of RNA Processing. By constructing explicit counterexamples, we see that the usual hypotheses cannot be weakened without loss of the theorem's conclusion.

\vspace{0.5cm}

\textbf{Solution 22:}

Solution approach: First recall the definitions and key theorems related to Translation from 7.28. Then apply the appropriate techniques to construct a rigorous argument. The solution must be self-contained and reference only the material from the course.

\vspace{0.5cm}

\textbf{Solution 23:}

The proof follows from the definitions and properties established in 7.28. The key insight is to apply the relevant theorem and verify the hypotheses hold. Detailed solution depends on the specific formulation of the theorem.

\vspace{0.5cm}

\textbf{Solution 24:}

The proof follows from the definitions and properties established in 7.28. The key insight is to apply the relevant theorem and verify the hypotheses hold. Detailed solution depends on the specific formulation of the theorem.

\vspace{0.5cm}

\textbf{Solution 25:}

The proof follows from the definitions and properties established in 7.28. The key insight is to apply the relevant theorem and verify the hypotheses hold. Detailed solution depends on the specific formulation of the theorem.

\vspace{0.5cm}

\end{document}